\documentclass[11pt]{article}
\usepackage{amsmath, amssymb, booktabs, geometry, hyperref, array, xcolor}
\geometry{margin=1in}
\hypersetup{colorlinks=true, linkcolor=blue, urlcolor=blue}
\setlength{\parskip}{0.5em}
\newcolumntype{L}[1]{>{\raggedright\arraybackslash}p{#1}}

\title{Simulation Study: Base Model and Sensitivity Scenarios\\
\large A Bayesian Framework for Causal Inference with Error-Prone Environmental Exposures}
\author{Johnson Vo}
\date{\today}

\begin{document}
\maketitle

\section{Motivation}

The simulation reflects the two-level structure of the application: area-level count
health outcomes for participants nested within Forward Sortation Areas (FSA or the first
three characters of a Canadian postal code), with area-level exposure assigned from a
spatiotemporal prediction model.

\section{Notation}

\begin{itemize}
\item $i = 1, \ldots, I$: area index (FSAs)
\item $A_i$: latent true area-level PM$_{2.5}$ exposure, unobserved
\item $Z_i$: observed error-prone surrogate from a spatiotemporal prediction model
(e.g. CanOSSEM)
\item $\mathbf{X}_i = (X_{i1}, \ldots, X_{i4})^\top$: area-level confounders affecting
both exposure and outcome
\item $D_i$: area-level calibration covariate explaining prediction error magnitude in
$Z_i$ (e.g. distance to nearest monitoring station, area deprivation)
\item $Y_i$: area-level count health outcome (number of incident depression cases)
\item $N_i$: area population size, used as Poisson offset
\item $\phi_i$: spatial random effect on the latent exposure surface, following the
CAR prior in equation~\eqref{eq:car}
\item $\mathbf{W}$: $I \times I$ binary adjacency matrix ($W_{ij} = 1$ if areas $i$ and
$j$ share a border)
\item $\mathbf{D}_W = \mathrm{diag}(\mathbf{W}\mathbf{1})$: diagonal degree matrix
(number of neighbours per area)
\item $\rho_A \in [0, 1)$: spatial correlation in the exposure CAR prior
\item $\sigma_\phi^2$: marginal variance of the spatial random effect $\phi_i$
\end{itemize}

As in the thesis proposal, $\mathbf{X}_i$ controls confounding while $D_i$ indexes the
measurement error mechanism. $D_i$ may overlap with $\mathbf{X}_i$ but is conceptually
distinct.

\section{Key Assumptions}

The simulation satisfies the following identification assumptions from the thesis proposal.

\begin{itemize}
\item \textbf{Non-differential surrogate:} $Y_i \perp Z_i \mid A_i, \mathbf{X}_i, D_i$.
Once we condition on the true exposure and covariates, the surrogate carries no additional
information about the outcome. Satisfied by construction since $Y_i$ is generated from
$A_i$ only and $Z_i$ never appears in the outcome equation.

\item \textbf{Transportability:} The error relationship is assumed transportable to the
target population. When direct calibration data are unavailable, inference about $A_i$ is
more model-dependent; this is examined in sensitivity scenario S2.
\end{itemize}

\section{Base Model}
\label{sec:base}

\subsection{Spatial Structure and CAR Prior}

Areas are arranged on a $\sqrt{I} \times \sqrt{I}$ synthetic grid with rook contiguity
(shared edges only). $I = 200$ balances the sample sizes used by Josey et al. (2023)
($n \in \{400, 800\}$) against the computational cost of BART.

Following thesis proposal Section 4.3, spatial dependence enters \emph{only} through the
latent exposure model via $\phi_i$. The spatial random effect follows the Conditional
Autoregressive (CAR) prior of Besag (1974):
\begin{equation}
\boldsymbol{\phi} \mid \rho_A, \sigma_\phi^2 \sim \mathcal{N}\!\left(\mathbf{0},\;
\sigma_\phi^2(\mathbf{D}_W - \rho_A\mathbf{W})^{-1}\right). \label{eq:car}
\end{equation}

\begin{itemize}
\item $\phi_i$ is a spatially structured random effect: neighbouring areas have similar
$\phi_i$ values, inducing correlation in their true exposures $A_i$
\item The precision matrix $(\mathbf{D}_W - \rho_A\mathbf{W})$ is sparse with adjacent
pairs having off-diagonal entry $-\rho_A$ and non-adjacent pairs having zero
\item $\rho_A = 0$: areas are spatially independent;
$\rho_A \to 1$: strong smoothing across neighbours
\item $\sigma_\phi^2 = 0.3$ is a design choice giving moderate spatial variation in
the exposure surface; in the application this parameter would be estimated from data
\item Spatial correlation in the outcome is induced indirectly through $A_i$; no
separate outcome random effect is added
\end{itemize}

\subsection{Covariates and Population Sizes}

\begin{itemize}
\item \textbf{Area confounders:} $X_{ik} \overset{\text{iid}}{\sim} \mathcal{N}(0, 1)$
for $k = 1, 2, 3, 4$
\item \textbf{Population sizes:} $\log(N_i) \sim \mathcal{N}(8.5, 1.2^2)$, giving
median $\approx 4{,}900$, 5th percentile $\approx 600$, 95th percentile $\approx 40{,}000$.
Motivated by Statistics Canada 2021 Census: Ontario has $\approx 500$ FSAs with total
population 14,223,942, ranging from dense urban FSAs in Toronto to sparse rural FSAs in
Northern Ontario.
\end{itemize}

\subsection{Calibration Covariate $D_i$}

$D_i$ captures area characteristics that determine prediction error in $Z_i$.

\paragraph{Base model: scalar $D_i$.}
$D_i \overset{\text{iid}}{\sim} \text{Gamma}(2, 1)$ (mean 2, variance 2), representing
distance to the nearest monitoring station.
\begin{itemize}
\item \textbf{Why Gamma?} Distance is strictly positive and right-skewed. Most Ontario
FSAs are moderately close to one of the 34 NAPS stations, but some rural areas are very
far. Shape 2 gives a unimodal distribution; Exponential (shape 1) would over-represent
very small values.
\item \textbf{Scientific motivation:} Gryparis et al. (2009) and Szpiro et al. (2011)
show that kriging prediction error variance is primarily determined by distance to the
monitor network. MELONS (Evangelopoulos et al., 2026) identified distance, deprivation,
and housing type as empirical determinants of PM$_{2.5}$ measurement error in London, UK;
we use their qualitative finding but not their specific values since they are not directly
transferable to Ontario.
\item \textbf{Ideally:} $D_i$ would be the actual Euclidean distance from each Ontario
FSA centroid to the 34 NAPS station coordinates, replacing the Gamma assumption with an
empirical distribution.
\end{itemize}

\paragraph{Extension: vector $D_i$.}
In a full implementation $D_i$ would be a vector with components such as:
$D_{i1} \sim \text{Gamma}(2,1)$ (distance to monitor),
$D_{i2} \sim \mathcal{N}(0,1)$ (income and employment indicators as a standardised
score), $D_{i3} \sim \text{Beta}(2,2)$ (proportion detached housing).
The error model would be $\omega^2(\mathbf{D}_i) = \exp(\eta_0 + \eta_1 D_{i1} +
\eta_2 D_{i3})$. Left as a future extension.

\subsection{Latent True Exposure}

The latent true exposure is generated from:
\begin{equation}
A_i \mid \mathbf{X}_i, \phi_i \sim \mathcal{N}\!\left(7 + 1.5X_{i1} - 1.5X_{i2} -
1.5X_{i3} + 1.5X_{i4} + \phi_i,\ 9\right), \label{eq:exposure}
\end{equation}
with $\sigma_A = 3$, where $\phi_i$ is the spatial random effect following the CAR prior
defined in equation~\eqref{eq:car} with parameters $(\rho_A, \sigma_\phi^2)$. The spatial
random effect $\phi_i$ induces correlation in $A_i$ across neighbouring areas: areas with
$\rho_A$ close to 1 have exposures that vary smoothly across the grid, reflecting a
regional wildfire plume, while $\rho_A = 0$ gives spatially independent exposures.

\begin{itemize}
\item \textbf{Mean = 7 $\mu$g/m$^3$:} Empirical annual average PM$_{2.5}$ in Ontario
(mean $\approx 7$, median $\approx 5.4$ $\mu$g/m$^3$)
\item \textbf{$\sigma_A = 3$:} Reflects realistic spatial variation
\item \textbf{Covariate coefficients $\pm 1.5$:} Provides meaningful confounding given
the exposure distribution
\item \textbf{Note on skewness:} Ontario PM$_{2.5}$ is heavily right-skewed (skewness
$\approx 9.2$, 99th percentile $\approx 26.4$, extreme values $> 200$ $\mu$g/m$^3$
during severe wildfire events). The Normal captures typical days but not extreme events.
\end{itemize}

\subsection{Measurement Error Model}

Two parameterisations of the measurement error model are considered. Both generate $Z_i$
as a function of the true exposure $A_i$ and calibration covariate $D_i$, and both
satisfy the non-differential assumption by construction since $Z_i$ does not appear in
the outcome equation.

In both versions the measurement errors $\varepsilon_i = Z_i - \mathbb{E}[Z_i \mid A_i, D_i]$
are assumed conditionally independent across areas given $(A_i, D_i)$. This is a
simplifying assumption: Gryparis et al.\ (2009) show that prediction errors from spatial
models are spatially correlated because nearby areas share the same monitoring station
data when their surrogates are generated. Incorporating spatially correlated errors is
listed as a potential extension in Section~\ref{sec:extensions}.

\paragraph{Version A: Classical error with Di-dependent mean.}

Following Wu et al.\ (2019), $D_i$ enters through the \emph{mean} of the surrogate,
introducing systematic directional bias:
\begin{equation}
Z_i \mid A_i, D_i \sim \mathcal{N}(\beta_{Z1} A_i + \beta_{Z2} D_i,\ \sigma_Z^2).
\label{eq:errorA}
\end{equation}

\begin{itemize}
\item \textbf{$\beta_{Z1}$}: attenuation. $\beta_{Z1} = 1$ centres the surrogate on the
truth; $\beta_{Z1} < 1$ attenuates high exposures. Wu et al.\ (2019) use $\beta_1 = 0.8$
as their default.
\item \textbf{$\beta_{Z2}$}: directional bias from $D_i$. When $\beta_{Z2} > 0$ the
surrogate overestimates true exposure in high-$D_i$ areas (far from monitors), consistent
with the MELONS finding of systematic overestimation in deprived and remote areas. When
$\beta_{Z2} = 0$ this reduces to standard classical error.
\item \textbf{$\sigma_Z^2$}: fixed homoscedastic noise. Heterogeneity across areas is
captured entirely through the mean shift $\beta_{Z2} D_i$.
\end{itemize}

Wu et al.\ (2019) write this in the regression calibration direction as
$E(X \mid W, D) = \gamma_0 + \gamma_1 W + \gamma_2^\top D$.

\paragraph{Version B: Classical error with Di-dependent variance.}

Alternatively, $D_i$ enters through the variance of the surrogate, relaxing the
homoscedasticity assumption of Josey et al.\ (2023) Assumption 2:
\begin{equation}
Z_i \mid A_i, D_i \sim \mathcal{N}(A_i,\ \sigma_Z^2(D_i)), \qquad
\sigma_Z^2(D_i) = \exp(\eta_0 + \eta_1 D_i). \label{eq:errorB}
\end{equation}

\begin{itemize}
\item The surrogate is unbiased on average ($E[Z_i \mid A_i] = A_i$) but noisier in
high-$D_i$ areas
\item $\eta_1 = 0$ recovers Josey's homoscedastic Assumption 2 exactly
\item $\exp(\eta_0 + \eta_1 D_i)$ ensures positivity and implies a log-linear
(proportional) relationship: a unit increase in $D_i$ multiplies the error variance by
$\exp(\eta_1)$, natural for distance-based quantities (Gryparis et al., 2009)
\item Unlike Version A, there is no attenuation or directional bias; the error is
purely a noise term whose magnitude depends on $D_i$
\end{itemize}

The two versions make different assumptions about how $D_i$ influences the measurement
error: Version A captures directional bias (the surrogate is systematically wrong in a
direction that depends on $D_i$) while Version B captures variance heterogeneity (the
surrogate is unbiased but less reliable in high-$D_i$ areas). Both are scientifically
motivated and are presented to the supervisor for a decision on which best reflects the
CanOSSEM error structure.

\subsection{Outcome Model}

The area-level count outcome follows a Poisson model with a log-linear predictor and
population offset, matching Josey et al. (2023):
\begin{equation}
Y_i \mid A_i, \mathbf{X}_i, N_i \sim \text{Poisson}(N_i \exp(\eta_i)),
\label{eq:outcome}
\end{equation}
\begin{equation}
\eta_i = -3 + f(A_i) - 0.5X_{i1} - 0.25X_{i2} + 0.25X_{i3} + 0.5X_{i4}.
\end{equation}

\begin{itemize}
\item \textbf{Population offset $\log(N_i)$:} $N_i$ scales the expected count so that
$\exp(\eta_i)$ is the event rate per person. A larger area will have more events simply
because it has more people; the offset removes this confound. In R this enters as an
additive term $\log(N_i)$ in the linear predictor or via the \texttt{offset} argument.
\item \textbf{Intercept $-3$:} $\exp(-3) \approx 5\%$ baseline event rate, representing
baseline depression incidence in the CLSA population (Raina et al., 2019)
\item \textbf{Confounder coefficients $\pm(0.5, 0.25)$:} Follow Josey et al. (2023)
\item \textbf{ERF $f(A_i)$:} Exposure-response function specified in sensitivity
scenario S4; three functional forms are evaluated
\item \textbf{No spatial random effect:} Consistent with thesis proposal Section 4.3;
spatial correlation in outcomes is induced indirectly through the spatially correlated
$A_i$ surface
\end{itemize}

$Z_i$ does not appear in equation~\eqref{eq:outcome}, satisfying the non-differential
assumption by construction.

\subsection{Causal Estimand}

The target estimand is the average potential outcome $\mu(a) = \mathbb{E}[Y_i(a)/N_i]$
(average event rate under intervention $A_i = a$ for all areas) and the causal contrast
$\Delta(a_1, a_0) = \mu(a_1) - \mu(a_0)$.

\begin{itemize}
\item The true $\mu(a)$ is computed by Monte Carlo over $N_\text{ref} = 100{,}000$
fixed covariate vectors drawn once, held constant across all replicates
\item Evaluation grid: 201 equally spaced values covering the bulk of the $A_i$
distribution, following Josey et al. (2023)
\item A pre-specified pairwise contrast (e.g. $a_1 = 15$ vs $a_0 = 5$ $\mu$g/m$^3$)
is also reported as a single summary metric
\end{itemize}

\section{Simulation Parameter Grid}

\begin{table}[h!]
\centering
\small
\caption{Simulation parameters.
Version A and Version B refer to the two error model parameterisations in
Section~\ref{sec:base}.}
\label{tab:params}
\begin{tabular}{L{3.5cm} L{1.5cm} L{2.8cm} L{5.4cm}}
\toprule
Parameter & Notation & Base / Range & Justification \\
\midrule
\multicolumn{4}{l}{\textit{Design}} \\
Number of areas & $I$ & 200 & Josey uses 400--800; reduced for BART cost \\
Replicates & $R$ & 500 & Josey et al. (2023) \\
ERF evaluation points & & 201 & Josey et al. (2023) \\
\midrule
\multicolumn{4}{l}{\textit{Covariates}} \\
Area confounders & $X_{ik}$ & $\mathcal{N}(0,1)$ & Josey et al. (2023) \\
Population size & $\log(N_i)$ & $\mathcal{N}(8.5, 1.44)$ & Statistics Canada 2021 Census \\
Calibration covariate & $D_i$ & Gamma$(2,1)$ & Monitor distance proxy;
Gryparis et al. (2009) \\
\midrule
\multicolumn{4}{l}{\textit{True exposure}} \\
Mean & $\mu_{A,i}$ & $7 + 1.5\mathbf{X}_i$ & Ontario PM$_{2.5}$ mean $\approx 7$
$\mu$g/m$^3$ \\
SD & $\sigma_A$ & 3 & Realistic spatial variation \\
Spatial correlation & $\rho_A$ & $\{0, 0.5, 0.9\}$ & 0 = Josey; 0.9 = strong
spatial smoothing \\
Spatial variance & $\sigma_\phi^2$ & 0.3 & Design choice \\
\midrule
\multicolumn{4}{l}{\textit{Error model Version A (Classical Error)}} \\
Attenuation & $\beta_{Z1}$ & $\{0.6, 0.8, 1.0\}$ & 1.0 = no attenuation;
0.8 = Wu et al. (2019) \\
Directional bias & $\beta_{Z2}$ & $\{-1, 0, 1, 2\}$ & 0 = no bias;
$>0$ = overestimation (MELONS) \\
Noise variance & $\sigma_Z^2$ & $\{1, 4, 9\}$ & Low/moderate/high noise \\
\midrule
\multicolumn{4}{l}{\textit{Error model Version B (variance heteroscedasticity)}} \\
Baseline variance & $\eta_0$ & $\{-1, 0, 1\}$ & Error at $D_i = 0$ \\
Di-slope & $\eta_1$ & $\{0, 0.5, 1.0, 1.5, 2.0\}$ & $\eta_1 = 0$ recovers
Josey Assumption 2 \\
\midrule
\multicolumn{4}{l}{\textit{Outcome}} \\
Intercept & & $-3$ & $\approx 5\%$ baseline event rate (CLSA) \\
ERF & $f(A_i)$ & Linear/Quadratic/Nonlinear & See S4 \\
Confounder coefs & $\boldsymbol{\gamma}$ & $(-0.5,-0.25, 0.25, 0.5)$ &
Josey et al. (2023) \\
\midrule
\multicolumn{4}{l}{\textit{Sensitivity scenarios}} \\
$D_i$ distribution & & Gamma$(2,0.5)$; Exp$(1)$ & Wider spread; Southern Ontario \\
$N_i$ distribution & & Uniform$(10,1000)$ & Josey comparability \\
Exposure distribution & & Log-normal & Extreme wildfire events \\
Number of areas & $I$ & $\{100, 200, 400\}$ & Scaling \\
\bottomrule
\end{tabular}
\end{table}

\section{Sensitivity Scenarios}
\label{sec:sensitivity}

\paragraph{S1: $D_i$ distribution.}
\textbf{Gamma(2, 0.5)}: mean 4, more heterogeneity between areas.
\textbf{Exponential(1)}: monotone decreasing, consistent with Southern Ontario station
concentration.
\textbf{Empirical}: actual FSA-to-NAPS distances from Statistics Canada shapefiles.

\paragraph{S2: Error model alternatives.}
\textbf{Version A vs Version B}: The two parameterisations in Section~\ref{sec:base} are
compared directly to assess whether directional bias (Version A) or variance
heteroscedasticity (Version B) better characterises CanOSSEM error structure.

\paragraph{S3: Spatial correlation.}
\textbf{Scenario B ($\rho_A = 0$)}: spatially independent exposures. Directly comparable
to Josey et al. (2023) which sets $\phi_i = 0$ for all areas.
\textbf{Scenario A ($\rho_A = 0.9$)}: high spatial smoothing, reflecting a sustained
regional wildfire event where nearby FSAs share similar smoke concentrations.

\paragraph{S4: Exposure-response functional forms.}
Three forms spanning simple to complex, crossed with spatial scenarios (S3) to give six
main simulation cells.

\textbf{Linear}: $f(A_i) = 0.10(A_i - 7)$.
Monotone increasing with no curvature. Coefficient 0.10 rescaled from Josey's 0.25 for
the Ontario exposure scale. The workhorse of air pollution epidemiology (Pope et al.,
2002).

\textbf{Quadratic}: $f(A_i) = 0.15(A_i - 7) - 0.008(A_i - 7)^2$.
Concave shape that rises steeply then flattens at high concentrations. Biologically
motivated by potential saturation of health pathways at high PM$_{2.5}$.

\textbf{Nonlinear (Josey cosine form)}: $f(A_i) = 0.10(A_i - 7) - 0.75\cos(\pi(A_i
- 5)/8) - 0.25(A_i - 7)X_{i1}$.
S-shaped curve with an exposure-covariate interaction. The interaction means the ERF
shape differs across areas depending on $X_{i1}$, which BART can capture but parametric
GLMs cannot. Primary misspecification stress test.

\paragraph{S5: Log-normal exposure.}
$\log(A_i) \sim \mathcal{N}(1.68, 0.54)$ matching Ontario moments: median 5.4, mean 7,
SD 5.9. Captures extreme wildfire events (99th percentile $\approx 26$ $\mu$g/m$^3$)
that the Normal cannot represent. ERF evaluation grid adapts to the 5th--95th percentile.

\paragraph{S6: Number of areas.}
$I \in \{100, 200, 400\}$. All must be perfect squares for the synthetic grid.

\section{Performance Metrics}

For each scenario and method (naive, oracle, proposed, regression calibration):

\begin{itemize}
\item \textbf{Bias}: $\hat{\mu}(a) - \mu(a)$ averaged over evaluation points and
replicates
\item \textbf{RMSE}: $\sqrt{R^{-1}\sum_r (\hat{\mu}_r(a) - \mu(a))^2}$
\item \textbf{Coverage}: fraction of 95\% credible intervals containing $\mu(a)$
\item \textbf{Interval width}: mean width of 95\% credible intervals
\item \textbf{Causal contrast bias}: $\hat{\Delta}(a_1, a_0) - \Delta(a_1, a_0)$ for
a pre-specified pair
\end{itemize}

Pointwise results at $a = 11$ $\mu$g/m$^3$ are also reported following Josey et al.
(2023) Table 1.

\section{Potential Extensions}
\label{sec:extensions}

\begin{itemize}
\item \textbf{Spatially correlated measurement errors:} The current model assumes
conditional independence of measurement errors $\varepsilon_i = Z_i -
\mathbb{E}[Z_i \mid A_i, D_i]$ across areas which is a simplifying assumption that
Gryparis et al.\ (2009) show is violated in practice: prediction errors from spatial
models are correlated across nearby areas because those predictions are generated using
the same monitoring station data. A natural extension is to give the error term its own
spatial covariance structure, either by making $D_i$ spatially structured via a second
CAR prior (so that nearby areas have correlated error variances), or by adding an
explicit spatially correlated residual $\psi_i \sim \text{CAR}(\rho_\varepsilon,
\sigma_\varepsilon^2)$ to the error model.

\item \textbf{Individual-level binary outcome:} Replacing the area-level Poisson count
with individual-level Bernoulli outcomes $Y_{ij} \sim \text{Bernoulli}(p_{ij})$ and
individual covariates $L_{ij}$, matching the CLSA data structure. The Bayesian bootstrap
would draw Dirichlet weights over all $N = \sum_i n_i$ individuals rather than over areas.

\item \textbf{Spatial random effect on outcome ($\nu_i$)}: A second CAR prior capturing
residual spatial clustering in health outcomes (Englert et al., 2025). Left out of the
base model per thesis proposal Section 4.3.
\item \textbf{Negative binomial outcome:} The Poisson model assumes
$\text{Var}(Y_i) = \mathbb{E}[Y_i]$, which is often violated in practice when
area-level counts are overdispersed due to unmeasured heterogeneity across FSAs. A
negative binomial outcome relaxes this by introducing a dispersion parameter $\xi > 0$:
\begin{equation*}
Y_i \mid A_i, \mathbf{X}_i, N_i \sim \text{NegBin}(\mu_i, \xi), \qquad
\mu_i = N_i \exp(\eta_i),
\end{equation*}
where $\mathbb{E}[Y_i] = \mu_i$ and $\text{Var}(Y_i) = \mu_i + \mu_i^2/\xi$. As
$\xi \to \infty$ this recovers the Poisson; small $\xi$ gives strong overdispersion.
Equivalently this arises from a Poisson-Gamma hierarchy as in Englert et al.\ (2025):
$Y_i \mid \theta_i \sim \text{Poisson}(\theta_i)$ where $\theta_i \sim \text{Gamma}(\xi,
\mu_i/\xi)$, marginalising over $\theta_i$ to give the NegBin. The dispersion parameter
$\xi = 1$ is a natural starting point and is the value used by Englert et al.\ (2025) in
their simulation. This extension tests whether the proposed method remains well-calibrated
when the outcome model is more variable than the base Poisson.
 
\item \textbf{Zero-inflated outcome}: For rare outcomes or small populations where
structural zeros arise.
 
\item \textbf{Log-normal exposure model:} The base model specifies $A_i$ on the original
scale as a Normal with a floor at zero. An alternative motivated by the right-skewed
nature of PM$_{2.5}$ concentrations and by the suggestion that $\log(A_i)$ should be
regressed on covariates is to model the log-exposure directly:
\begin{equation*}
\log(A_i) \mid \mathbf{X}_i, \phi_i \sim \mathcal{N}\!\left(\gamma_0 + \gamma_1 X_{i1}
+ \gamma_2 X_{i2} + \gamma_3 X_{i3} + \gamma_4 X_{i4} + \phi_i,\
\sigma_{\log A}^2\right),
\end{equation*}
so that $A_i = \exp(\log A_i) > 0$ always. The parameters are matched to Ontario
PM$_{2.5}$ empirical moments: setting $\gamma_0 = 1.68$ and $\sigma_{\log A} = 0.73$
gives $\mathbb{E}[A_i] \approx 7$ $\mu$g/m$^3$ and $\text{SD}(A_i) \approx 5.9$
$\mu$g/m$^3$, matching the provincial monitoring data. The error model in Version A
would then operate on the log scale as $\log(Z_i) \mid \log(A_i), D_i \sim
\mathcal{N}(\beta_{Z1}\log(A_i) + \beta_{Z2}D_i, \sigma_Z^2)$, giving multiplicative
error on the original scale. The causal estimand $\mu(a)$ is computed at values of $a$
on the original scale, but the exposure-response curve is evaluated over a grid adapted
to the 5th--95th percentile of the log-normal distribution rather than a fixed interval.
\end{itemize}


\end{document}